% =====================================================================
\chapter{Combinaciones, números combinatorios y doble conteo}
\label{tema:2}
% =====================================================================

Este tema cuenta \emph{selecciones no ordenadas}: subconjuntos, multiconjuntos y repartos. Aparecen los números combinatorios, el binomio y el multinomio de Newton, la técnica del doble conteo y la forma general del principio de inclusión--exclusión. Se siguen las secciones §14.5 a §14.7 de Lehman, Leighton y Meyer~\cite{mcs}; las secciones §1.4 a §1.6 de Ross~\cite{ross}; las secciones §1.4 y §1.5, «Story proofs», de Blitzstein y Hwang~\cite{bh}, cuyo estilo de «demostración mediante una historia» se adopta aquí, y la sección~3.2 de Grinstead y Snell~\cite{gs}.

\section{Combinaciones y números combinatorios}

\begin{definicion}[Combinaciones]
Sea $A$ un conjunto con $n$ elementos y $0\le k\le n$. Una \emph{combinación} (sin repetición) de orden $k$ de $A$ es un subconjunto de $A$ con $k$ elementos. Su número se denota $\binom{n}{k}$ y se llama \emph{número combinatorio} («$n$ sobre $k$»). Por convenio, $\binom{n}{k}=0$ si $k<0$ o $k>n$.
\end{definicion}

\begin{teorema}\label{teo:t2-binom}
Para $0\le k\le n$,
\[
\binom{n}{k}=\frac{n^{\underline{k}}}{k!}=\frac{n!}{k!\,(n-k)!}.
\]
\end{teorema}

\begin{proof}
Sea $V$ el conjunto de variaciones sin repetición de orden $k$ de $A$ y $C$ el de subconjuntos de $k$ elementos. La aplicación $V\to C$ que olvida el orden, $(a_1,\dots,a_k)\mapsto\{a_1,\dots,a_k\}$, es sobreyectiva y cada subconjunto tiene exactamente $k!$ preimágenes, sus $k!$ ordenaciones. Por la regla de la división (\cref{prop:t1-division}), $\abs{V}=k!\,\abs{C}$, y $\abs{V}=n^{\underline{k}}$ por el \cref{teo:t1-variaciones}.
\end{proof}

\begin{proposicion}[Propiedades básicas]\label{prop:t2-props}
Para enteros $n\ge1$ y $k$:
\begin{enumerate}
  \item simetría: $\binom{n}{k}=\binom{n}{n-k}$;
  \item fórmula de Pascal: $\binom{n}{k}=\binom{n-1}{k}+\binom{n-1}{k-1}$;
  \item absorción: $k\binom{n}{k}=n\binom{n-1}{k-1}$.
\end{enumerate}
\end{proposicion}

\begin{proof}
Se dan demostraciones combinatorias; todas pueden comprobarse también con la fórmula del \cref{teo:t2-binom}.

(i) La aplicación $X\mapsto A\setminus X$ es una biyección entre los subconjuntos de $k$ elementos y los de $n-k$ elementos (es su propia inversa).

(ii) Se fija $a\in A$. Los subconjuntos de $k$ elementos de $A$ se dividen en dos clases disjuntas: los que no contienen a $a$, que son los subconjuntos de $k$ elementos de $A\setminus\{a\}$, en número $\binom{n-1}{k}$, y los que contienen a $a$, que se corresponden biyectivamente (quitando $a$) con los subconjuntos de $k-1$ elementos de $A\setminus\{a\}$, en número $\binom{n-1}{k-1}$. Se aplica la regla de la suma.

(iii) Se cuentan de dos formas los pares $(X,p)$ con $X\subseteq A$, $\abs{X}=k$ y $p\in X$ (un comité de $k$ personas con un presidente que pertenece a él). Eligiendo primero el comité y luego el presidente: $\binom{n}{k}\cdot k$. Eligiendo primero el presidente y luego los otros $k-1$ miembros entre las $n-1$ personas restantes: $n\binom{n-1}{k-1}$.
\end{proof}

La fórmula de Pascal permite construir el \emph{triángulo de Pascal}, en el que cada número es la suma de los dos que tiene encima:
\[
\begin{array}{c}
1\\
1\quad 1\\
1\quad 2\quad 1\\
1\quad 3\quad 3\quad 1\\
1\quad 4\quad 6\quad 4\quad 1\\
1\quad 5\quad 10\quad 10\quad 5\quad 1
\end{array}
\]

\section{Binomio y multinomio de Newton}

\begin{teorema}[Binomio de Newton]\label{teo:t2-binomio}
Para $x,y$ reales y $n\ge0$ entero,
\[
(x+y)^{n}=\sum_{k=0}^{n}\binom{n}{k}x^{k}y^{n-k}.
\]
\end{teorema}

\begin{proof}
Al desarrollar el producto $(x+y)(x+y)\cdots(x+y)$ de $n$ factores, sin conmutar, se obtiene la suma de las $2^{n}$ palabras que resultan de elegir $x$ o $y$ en cada factor. Una palabra con $k$ letras $x$ vale $x^{k}y^{n-k}$, y hay tantas como formas de elegir los $k$ factores de los que se toma $x$, es decir, $\binom{n}{k}$.
\end{proof}

El módulo~0.3 demuestra este teorema por inducción; la demostración anterior es la combinatoria, que explica \emph{por qué} aparecen los números combinatorios.

\begin{corolario}\label{cor:t2-sumas}
Para $n\ge 0$, $\sum_{k=0}^{n}\binom{n}{k}=2^{n}$, y para $n\ge1$, $\sum_{k=0}^{n}(-1)^{k}\binom{n}{k}=0$.
\end{corolario}

\begin{proof}
Basta tomar $x=y=1$ y $x=-1$, $y=1$ en el \cref{teo:t2-binomio}. La primera identidad tiene también una demostración combinatoria: ambos lados cuentan los subconjuntos de un conjunto de $n$ elementos, el izquierdo agrupándolos por tamaño~\cite[\emph{Strategic Practice}~1, problema~5]{stat110}.
\end{proof}

\begin{definicion}[Coeficiente multinomial]
Sean $n=k_1+\dots+k_r$ con $k_i\ge0$ enteros. Se define
\[
\binom{n}{k_1,k_2,\dots,k_r}=\frac{n!}{k_1!\,k_2!\cdots k_r!}.
\]
\end{definicion}

\begin{teorema}[Permutaciones con repetición]\label{teo:t2-multinomial}
El número de sucesiones de longitud $n$ formadas con $k_1$ copias del símbolo $s_1$, $k_2$ del símbolo $s_2$, \dots, $k_r$ del símbolo $s_r$ es $\binom{n}{k_1,\dots,k_r}$. Equivalentemente, es el número de formas de repartir $n$ objetos distintos en $r$ cajas distinguibles, con $k_i$ objetos en la caja $i$.
\end{teorema}

\begin{proof}
Una sucesión queda determinada eligiendo las $k_1$ posiciones de $s_1$, después las $k_2$ posiciones de $s_2$ entre las $n-k_1$ restantes, y así sucesivamente. Por la regla del producto generalizada,
\[
\binom{n}{k_1}\binom{n-k_1}{k_2}\cdots\binom{n-k_1-\dots-k_{r-1}}{k_r}
=\frac{n!}{k_1!\,(n-k_1)!}\cdot\frac{(n-k_1)!}{k_2!\,(n-k_1-k_2)!}\cdots=\frac{n!}{k_1!\cdots k_r!},
\]
porque el producto es telescópico. La segunda formulación es equivalente: la sucesión indica en qué caja va cada objeto.
\end{proof}

\begin{ejemplo}
La palabra \texttt{BANANA} tiene 6 letras: una \texttt{B}, tres \texttt{A} y dos \texttt{N}. Tiene $\binom{6}{1,3,2}=\frac{720}{1\cdot6\cdot2}=60$ anagramas distintos. Este recuento es el de la «regla del contable» de~\cite[§14.6.2]{mcs}.
\end{ejemplo}

\begin{teorema}[Multinomio de Newton]\label{teo:t2-multinomio}
Para reales $x_1,\dots,x_r$ y $n\ge0$,
\[
(x_1+\dots+x_r)^{n}=\sum_{\substack{k_1+\dots+k_r=n\\ k_i\ge0}}\binom{n}{k_1,\dots,k_r}x_1^{k_1}\cdots x_r^{k_r}.
\]
\end{teorema}

\begin{proof}
Como en el binomio: el desarrollo es la suma de las $r^{n}$ palabras de longitud $n$ sobre $\{x_1,\dots,x_r\}$, y las que contienen $k_i$ veces $x_i$ son $\binom{n}{k_1,\dots,k_r}$ por el \cref{teo:t2-multinomial}.
\end{proof}

\section{Combinaciones con repetición}

\begin{definicion}
Una \emph{combinación con repetición} de orden $k$ de un conjunto de $n$ elementos es una selección no ordenada de $k$ elementos en la que se permite repetir; equivale a un vector $(x_1,\dots,x_n)$ de enteros $x_i\ge 0$ con $x_1+\dots+x_n=k$, donde $x_i$ es el número de veces que se elige el elemento $i$.
\end{definicion}

\begin{teorema}[Estrellas y barras]\label{teo:t2-estrellas}
El número de soluciones enteras de $x_1+\dots+x_n=k$ con todos los $x_i\ge0$ es $\binom{n+k-1}{k}=\binom{n+k-1}{n-1}$. El número de soluciones con todos los $x_i\ge1$ es $\binom{k-1}{n-1}$ (para $k\ge n\ge 1$).
\end{teorema}

\begin{proof}
A cada solución $(x_1,\dots,x_n)$ se le asocia la cadena formada por $x_1$ estrellas, una barra, $x_2$ estrellas, una barra, \dots, y $x_n$ estrellas: tiene $k$ estrellas y $n-1$ barras. Recíprocamente, toda cadena de $k$ estrellas y $n-1$ barras determina una única solución (los $x_i$ son los tamaños de los bloques de estrellas separados por las barras). Es una biyección, y las cadenas se cuentan eligiendo las posiciones de las $n-1$ barras entre $n+k-1$: $\binom{n+k-1}{n-1}$. Para la segunda afirmación, el cambio $y_i=x_i-1$ transforma biyectivamente las soluciones con $x_i\ge1$ en las soluciones de $y_1+\dots+y_n=k-n$ con $y_i\ge0$, que son $\binom{k-1}{n-1}$.
\end{proof}

\begin{table}[h]
\centering
\begin{tabular}{@{}lcc@{}}
\toprule
Selección de $k$ entre $n$ & sin repetición & con repetición \\
\midrule
ordenada (importa el orden) & $n^{\underline{k}}=\dfrac{n!}{(n-k)!}$ & $n^{k}$ \\[8pt]
no ordenada & $\dbinom{n}{k}$ & $\dbinom{n+k-1}{k}$ \\
\bottomrule
\end{tabular}
\caption{Los cuatro modelos básicos de selección. El caso con repetición y sin orden es el único en que la fórmula no es «evidente».}
\label{tab:t2-modelos}
\end{table}

\section{Doble conteo y demostraciones biyectivas}

Una \emph{demostración por doble conteo} establece una identidad $a=b$ mostrando que ambos lados cuentan el mismo conjunto de dos maneras distintas; una \emph{demostración biyectiva} la establece construyendo una biyección entre un conjunto de $a$ elementos y otro de $b$ elementos. Ya se han usado ambas en la \cref{prop:t2-props}. El esquema conviene escribirlo siempre de forma explícita: (1) describir el conjunto que se cuenta (una «historia»); (2) contarlo de una forma, obteniendo el lado izquierdo; (3) contarlo de otra, obteniendo el lado derecho~\cite[§1.5]{bh}.

\begin{ejemplo}[Caminos en una cuadrícula]\label{ej:t2-caminos}
Un camino reticular de $(0,0)$ a $(a,b)$, con $a,b\ge 0$ enteros, es una sucesión de pasos unitarios hacia la derecha ($D$) o hacia arriba ($U$). Todo camino tiene $a$ pasos $D$ y $b$ pasos $U$, y queda determinado por la posición de los pasos $D$ entre los $a+b$; hay $\binom{a+b}{a}$. Clasificando los caminos de $(0,0)$ a $(a,b)$ según su último paso se obtiene de nuevo la fórmula de Pascal: $\binom{a+b}{a}=\binom{a+b-1}{a-1}+\binom{a+b-1}{a}$.
\end{ejemplo}

\begin{proposicion}[Identidad de Vandermonde]\label{prop:t2-vandermonde}
Para enteros $m,n,k\ge0$, $\displaystyle\binom{m+n}{k}=\sum_{j=0}^{k}\binom{m}{j}\binom{n}{k-j}$.
\end{proposicion}

La demostración se pide en el \cref{eje:2-vandermonde}: un grupo de $m$ mujeres y $n$ hombres, y comités de $k$ personas clasificados por el número de mujeres.

\section{Inclusión--exclusión general}

\begin{teorema}[Inclusión--exclusión]\label{teo:t2-ie}
Para conjuntos finitos $A_1,\dots,A_m$,
\[
\Bigl|\bigcup_{i=1}^{m}A_i\Bigr|=\sum_{j=1}^{m}(-1)^{j+1}\sum_{1\le i_1<\dots<i_j\le m}\abs{A_{i_1}\cap\dots\cap A_{i_j}}.
\]
\end{teorema}

\begin{proof}
Se comprueba que cada elemento $x$ de la unión contribuye exactamente~1 al lado derecho (los que no están en la unión contribuyen~0 a ambos lados). Sea $r\ge1$ el número de conjuntos $A_i$ que contienen a $x$. Una intersección de $j$ conjuntos contiene a $x$ si y solo si los $j$ conjuntos están entre esos $r$, lo que ocurre para $\binom{r}{j}$ de las intersecciones de $j$ conjuntos. La contribución de $x$ es, pues,
\[
\sum_{j=1}^{r}(-1)^{j+1}\binom{r}{j}=1-\sum_{j=0}^{r}(-1)^{j}\binom{r}{j}=1-0=1,
\]
por el \cref{cor:t2-sumas}.
\end{proof}

\begin{ejemplo}[Desarreglos]\label{ej:t2-desarreglos}
Un \emph{desarreglo} de $[n]$ es una permutación $\sigma$ sin puntos fijos: $\sigma(i)\neq i$ para todo $i$. Sea $A_i$ el conjunto de permutaciones con $\sigma(i)=i$. Una intersección de $j$ de ellos fija $j$ puntos y permuta libremente los demás, luego tiene $(n-j)!$ elementos. Por el \cref{teo:t2-ie}, el número de permutaciones con algún punto fijo es $\sum_{j=1}^{n}(-1)^{j+1}\binom{n}{j}(n-j)!$, y el de desarreglos es
\[
D_n=n!-\sum_{j=1}^{n}(-1)^{j+1}\frac{n!}{j!}=n!\sum_{j=0}^{n}\frac{(-1)^{j}}{j!}.
\]
Como $\sum_{j\ge0}(-1)^{j}/j!=e^{-1}$, la proporción de desarreglos $D_n/n!$ tiende a $1/e\approx0{,}368$. Por ejemplo, $D_4=24\bigl(1-1+\tfrac12-\tfrac16+\tfrac1{24}\bigr)=9$.
\end{ejemplo}

\begin{observacion}[Conexión con el curso]
En el módulo~22, la métrica \emph{pass@$k$} para evaluar código generado se estima a partir de $n$ muestras de las que $c$ son correctas: es la probabilidad de que un subconjunto aleatorio de $k$ muestras contenga alguna correcta, $1-\binom{n-c}{k}\big/\binom{n}{k}$, un cociente de números combinatorios~\cite[§2.1]{chen}. En el módulo~1, la verosimilitud de un corpus depende solo de los recuentos de cada token, y el número de secuencias con unos recuentos dados es un coeficiente multinomial. En el módulo~12, el algoritmo BPE cuenta pares de símbolos adyacentes: una secuencia de longitud $L$ tiene $L-1$ pares adyacentes, mientras que un conjunto de $K$ símbolos admite $K^{2}$ pares ordenados posibles.
\end{observacion}

\begin{observacion}[Errores frecuentes]
(1) Usar $\binom{n}{k}$ cuando importa el orden, o $n^{\underline{k}}$ cuando no importa. (2) Olvidar dividir por el número de ordenaciones de grupos \emph{indistinguibles}: repartir 12 personas en tres equipos de 4 sin nombre no es lo mismo que en los equipos A, B y C (\cref{eje:2-equipos}). (3) Contar «al menos un par» eligiendo el par y luego el resto libremente, lo que repite manos. (4) Confundir combinaciones con repetición con variaciones con repetición: el número de multiconjuntos de tamaño $k$ \emph{no} es $n^{k}/k!$. (5) En una demostración por doble conteo, no decir explícitamente qué conjunto se cuenta.
\end{observacion}

\begin{lecturas}
\item Lehman, Leighton y Meyer~\cite{mcs}, §14.5 «Counting Subsets», §14.6 «Sequences with Repetitions» y §14.7 «Counting Practice: Poker Hands».
\item Ross~\cite{ross}, §1.4 «Combinations», §1.5 «Multinomial Coefficients» y §1.6 «The Number of Integer Solutions of Equations».
\item Blitzstein y Hwang~\cite{bh}, §1.4 «How to count» y §1.5 «Story proofs».
\item Grinstead y Snell~\cite{gs}, §3.2 «Combinations».
\end{lecturas}

\section*{Ejercicios del tema 2}
\addcontentsline{toc}{section}{Ejercicios del tema 2}

\begin{ejercicio}{1}{eje:2-calcular}
Calcule $\binom{10}{3}$, $\binom{10}{7}$, $\binom{52}{5}$ y el coeficiente de $x^{3}$ en $(1+x)^{7}$. Explique sin hacer cuentas por qué $\binom{10}{3}=\binom{10}{7}$.
\end{ejercicio}

\begin{ejercicio}{1}{eje:2-modelo}
Identifique el modelo de la \cref{tab:t2-modelos} y cuente: (a) podios (oro, plata, bronce) entre 8 atletas; (b) comisiones de 3 personas entre 8; (c) secuencias de 3 tokens de un vocabulario de 8; (d) bolsas de 3 caramelos de 8 sabores, con sabores posiblemente repetidos.
\end{ejercicio}

\begin{ejercicio}{2}{eje:2-poker}
Una mano de póquer es un subconjunto de 5 cartas de una baraja francesa (13 valores, 4 palos). Cuente las manos con (a) póquer (cuatro cartas del mismo valor); (b) full (tres de un valor y dos de otro); (c) doble pareja (dos parejas de valores distintos y una quinta carta de un tercer valor). Exprese cada resultado como fracción del total de manos.
\fuente{adaptado de Lehman, Leighton y Meyer~\cite{mcs}, §14.7 «Counting Practice: Poker Hands».}
\end{ejercicio}

\begin{ejercicio}{2}{eje:2-equipos}
(a) ¿De cuántas formas pueden repartirse 12 personas en tres equipos de tamaños 2, 5 y 5 (los equipos no tienen nombre)? (b) ¿Y en tres equipos de 4?
\fuente{adaptado de Harvard Stat 110~\cite{stat110}, \emph{Homework}~1, problema~2.}
\end{ejercicio}

\begin{ejercicio}{2}{eje:2-anagramas}
Un anagrama de una palabra es una ordenación distinta de sus letras. Cuente los anagramas de \texttt{MISSISSIPPI} y los de \texttt{PROBABILIDAD}. ¿En cuántos anagramas de \texttt{MISSISSIPPI} aparecen las cuatro \texttt{S} juntas?
\end{ejercicio}

\begin{ejercicio}{2}{eje:2-coeficientes}
Halle el coeficiente de $x^{3}y^{5}$ en $(2x-y)^{8}$ y el de $x^{2}y^{3}z$ en $(x+y+z)^{6}$. ¿Cuántos términos distintos tiene el desarrollo de $(x+y+z)^{6}$?
\end{ejercicio}

\begin{ejercicio}{2}{eje:2-ecuaciones}
Cuente las soluciones enteras de $x_1+x_2+x_3+x_4=10$ (a) con $x_i\ge0$; (b) con $x_i\ge1$; (c) con $x_i\ge0$ y $x_1\ge2$.
\end{ejercicio}

\begin{ejercicio}{2}{eje:2-caminos}
Cuente los caminos reticulares (\cref{ej:t2-caminos}) (a) de $(0,0)$ a $(5,6)$; (b) de $(0,0)$ a $(5,6)$ que pasan por $(2,3)$; (c) de $(0,0)$ a $(5,6)$ que \emph{no} pasan por $(2,3)$.
\fuente{adaptado de Harvard Stat 110~\cite{stat110}, \emph{Strategic Practice}~1, problema~3 (con otros datos).}
\end{ejercicio}

\begin{ejercicio}{2}{eje:2-estaciones}
Se anota la estación del año (primavera, verano, otoño o invierno) en que nació cada una de 7 personas. ¿Cuántas de las $4^{7}$ asignaciones posibles de estaciones incluyen las cuatro estaciones? ¿Qué fracción del total representan?
\fuente{adaptado de Harvard Stat 110~\cite{stat110}, \emph{Strategic Practice}~2, sección «Inclusion-Exclusion», problema~1.}
\end{ejercicio}

\begin{ejercicio}{2}{eje:2-desarreglos}
Cinco amigos organizan un amigo invisible: cada uno saca al azar el nombre de otro de una bolsa con los cinco nombres. (a) ¿Cuántos sorteos hay en los que nadie se saca a sí mismo? (b) ¿Qué fracción de los $5!$ sorteos representan? Compárela con $1/e$.
\end{ejercicio}

\begin{ejercicio}{2}{eje:2-ositos}
Una bolsa de gominolas contiene entre 30 y 50 ositos, cada uno de uno de 5 sabores. Dos bolsas se consideran iguales si tienen el mismo número de ositos de cada sabor. ¿Cuántas bolsas distintas hay? Exprese la respuesta como diferencia de dos números combinatorios, no como una suma larga.
\fuente{adaptado de Harvard Stat 110~\cite{stat110}, \emph{Homework}~1, problema~7(b).}
\end{ejercicio}

\begin{ejercicio}{3}{eje:2-palo}
Demuestre mediante una historia (doble conteo) que, para enteros $n\ge k\ge0$,
\[
\binom{k}{k}+\binom{k+1}{k}+\dots+\binom{n}{k}=\binom{n+1}{k+1}.
\]
\emph{Indicación:} ordene $n+1$ personas por edad y clasifique los grupos de $k+1$ según quién es el mayor del grupo.
\fuente{adaptado de Harvard Stat 110~\cite{stat110}, \emph{Homework}~1, problema~7(a).}
\end{ejercicio}

\begin{ejercicio}{3}{eje:2-sumak}
(a) Usando la absorción (\cref{prop:t2-props}), demuestre que $\sum_{k=0}^{n}k\binom{n}{k}=n2^{n-1}$ para $n\ge1$. (b) Demuestre que, para $n\ge2$, $\sum_{k=2}^{n}k(k-1)\binom{n}{k}=n(n-1)2^{n-2}$, contando de dos maneras los comités con un presidente y un vicepresidente distintos.
\fuente{el apartado (b) está adaptado de MIT 6.041~\cite{mit6041f10}, \emph{Problem Set}~3, problema~5.}
\end{ejercicio}

\begin{ejercicio}{3}{eje:2-vandermonde}
Demuestre la identidad de Vandermonde (\cref{prop:t2-vandermonde}) por doble conteo y deduzca que $\sum_{j=0}^{n}\binom{n}{j}^{2}=\binom{2n}{n}$.
\end{ejercicio}

\begin{ejercicio}{3}{eje:2-pares}
Sea $A$ un conjunto con $n\ge1$ elementos y $a\in A$ fijo. (a) Demuestre que la aplicación $X\mapsto X\dsim\{a\}$ (añadir $a$ si no está, quitarlo si está) es una biyección entre los subconjuntos de $A$ de cardinal par y los de cardinal impar. (b) Deduzca, sin usar el binomio, que $\sum_{k}(-1)^{k}\binom{n}{k}=0$ y que $A$ tiene $2^{n-1}$ subconjuntos de cardinal par.
\end{ejercicio}
